% !TEX program = pdflatex
\documentclass[11pt,a4paper]{article}

\usepackage[utf8]{inputenc}
\usepackage[T1]{fontenc}
\usepackage[spanish,es-nodecimaldot]{babel}
\usepackage{amsmath,amssymb,amsthm,mathtools,mathrsfs}
\usepackage{geometry}
\usepackage{enumitem}
\usepackage[expansion=false]{microtype}
\usepackage{hyperref}
\usepackage{cleveref}
\usepackage{xcolor}
\geometry{margin=2.7cm}
\setlength{\parindent}{0pt}
\setlength{\parskip}{0.6em}

\hypersetup{
	colorlinks=true,
	linkcolor=blue!60!black,
	urlcolor=blue!60!black,
	citecolor=blue!60!black
}

%------------------------------------------------
% Entornos
%------------------------------------------------
\theoremstyle{definition}
\newtheorem{definicion}{Definición}[section]
\newtheorem{ejemplo}{Ejemplo}[section]

\theoremstyle{plain}
\newtheorem{proposicion}{Proposición}[section]
\newtheorem{teorema}{Teorema}[section]
\newtheorem{corolario}{Corolario}[section]
\newtheorem{lema}{Lema}[section]

\theoremstyle{remark}
\newtheorem{observacion}{Observación}[section]

%------------------------------------------------
% Comandos
%------------------------------------------------
\newcommand{\C}{\mathbb C}
\newcommand{\R}{\mathbb R}
\newcommand{\Z}{\mathbb Z}
\newcommand{\N}{\mathbb N}
\newcommand{\I}{[0,1]}
\newcommand{\Hol}{\mathcal H}
\newcommand{\OO}{\mathcal O}
\newcommand{\id}{\operatorname{id}}
\newcommand{\restr}[2]{\left.#1\right|_{#2}}

\title{\textbf{Notas de EDO}}
\author{}
\date{}

\begin{document}
	\maketitle
	\textbf{Objetivo}
	Notas del curso auto suficientes para los temas abarcados semana a semana
	\section{Soluciones Maximales}
	Sea $U \subseteq \mathbb{R}^n$ abierto, $x_0 \in U$, $t_0 \in \mathbb{R}$ y $X \in \mathcal{X}^1(U)$. Por el Teorema de Picard, el P.V.I.
	\begin{equation*}
		\left| 
		\begin{array}{lcl}
			x' &=& X(x) \\
			x(t_0) &=& x_0
		\end{array}
		\right.
	\end{equation*}
	admite una única solución local $\varphi : I_\alpha[t_0] \to B_r[x_0] \subseteq U$ donde $\alpha = \frac{r}{M}$ y siendo $M = \|X\|_{C(B_r[x_0])}$.
	
	Una pregunta natural sería: ¿Podemos extender la solución dada a un intervalo $J$ que contenga a $I_\alpha[t_0]$ y tal que la extensión siga siendo solución del P.V.I.? Como veremos a continuación, la respuesta a esta interrogante es afirmativa.
	
	\begin{definicion}
		Sean $U \subseteq \mathbb{R}^n$ abierto, $X \in \mathcal{X}^0(U)$, $x_0 \in U$, $t_0 \in \mathbb{R}$ y consideremos el P.V.I.:
		\begin{equation}\label{eq:517}
			\left| 
			\begin{array}{lcl}
				x' &=& X(x) \\
				x(t_0) &=& x_0
			\end{array}
			\right.
		\end{equation}
		Una solución $\varphi_M : I_M \to \mathbb{R}^n$ de \eqref{eq:517} es llamada \textit{solución maximal} y su dominio $I_M$ es llamado \textit{intervalo maximal} si y sólo si toda solución $\psi : J \to \mathbb{R}^n$ del P.V.I. \eqref{eq:517} tal que $I_M \subseteq J$ y $\left.\psi\right|_{I_M} = \varphi_M$ implica que $I_M = J$.
	\end{definicion}
	\noindent \textbf{Observaciones:}
	\begin{enumerate}
		\item Una solución $\varphi_M$ de \eqref{eq:517} es maximal si y sólo si no admite extensiones las cuales sean también solución del P.V.I. \eqref{eq:517}.
		\item Tanto la solución maximal $\varphi_M$ como el intervalo maximal $I_M$ dependen del par $(t_0, x_0) \in \mathbb{R} \times U$. Por esta razón denotaremos $\varphi_M = \varphi_{(t_0, x_0)}$ a la solución maximal y por $I_M = I(t_0, x_0)$ a su intervalo maximal.
		\item Si $X \in \mathcal{X}^0(U)$ entonces su P.V.I. asociado puede admitir más de una solución maximal. En efecto, consideremos el P.V.I.
		\begin{equation*}
			\left| 
			\begin{array}{lcl}
				x' &=& 5x^{4/5} \\
				x(0) &=& 0
			\end{array}
			\right.
		\end{equation*}
		entonces $\{\varphi_c\}_{c > 0}$ donde
		\[
		\begin{array}{rcccl}
			\varphi_c : & \mathbb{R} &\to& \mathbb{R} \\
			& t &\mapsto& \varphi(t) = \begin{cases} (t - c)^5, & \text{si } t \ge c \\ 0, & \text{si } t \le c \end{cases}
		\end{array}
		\]
		es una familia infinita (no numerable) de soluciones maximales del P.V.I. dado.
	\end{enumerate}
	
	El siguiente resultado establece que toda solución de un P.V.I. definida en un intervalo cerrado puede ser extendida a un intervalo abierto.
	
	\noindent \textbf{Proposición 5.4.1} Sea $U \subseteq \mathbb{R}^n$ abierto, $X \in \mathcal{X}^0(U)$ y suponga que para cualquier $(t_0, x_0) \in \mathbb{R} \times U$ el P.V.I.
	\begin{equation}\label{eq:518}
		\left| 
		\begin{array}{lcl}
			x' &=& X(x) \\
			x(t_0) &=& x_0
		\end{array}
		\right.
	\end{equation}
	admite una única solución. Si $\varphi : [a,b] \to U$ (donde $t_0 \in [a,b]$) es solución del P.V.I. \eqref{eq:518} entonces ella puede ser extendida a un intervalo abierto $]a - \delta, b + \epsilon[$.
	
	\noindent \textbf{Demostración.} Sea $x_1 = \varphi(b) \in U$ y consideremos el P.V.I.
	\begin{equation}\label{eq:519}
		\left| 
		\begin{array}{lcl}
			x' &=& X(x) \\
			x(b) &=& x_1
		\end{array}
		\right.
	\end{equation}
	
	Por Picard, existe $\phi : I_\epsilon[b] \to U$ solución de \eqref{eq:519}. Defino $\psi : [a, b + \epsilon[ \to U$ como
	\[
	\psi(t) = \begin{cases} 
		\varphi(t), & a \le t \le b \\ 
		\phi(t), & b < t < b + \epsilon 
	\end{cases}
	\]
	
	Afirmo que $\psi$ es diferenciable en $b$. En efecto:
	\[
	\lim_{t \to b^+} \frac{\psi(t) - \psi(b)}{t - b} = \lim_{t \to b^+} \frac{\phi(t) - \phi(b)}{t - b} = \phi'_+(b) = \phi'(b) = X(\phi(b)) = X(x_1)
	\]
	Análogamente:
	\[
	\lim_{t \to b^-} \frac{\psi(t) - \psi(b)}{t - b} = \lim_{t \to b^-} \frac{\varphi(t) - \varphi(b)}{t - b} = \varphi'_-(b) = \varphi'(b) = X(\varphi(b)) = X(x_1)
	\]
	
	Esto prueba la afirmación. Por otro lado, es claro que $\psi$ es solución del P.V.I. \eqref{eq:518} y por tanto, es una extensión de $\varphi$. Se procede de manera análoga con el extremo inferior. \hfill $\square$
	
	\noindent \textbf{Observación:} De ahora en adelante podemos suponer que las soluciones locales de un P.V.I. están definidas en un intervalo abierto.
	
	A continuación, probaremos que si el P.V.I. \eqref{eq:517} admite una única solución local para cada elección de $t_0 \in \mathbb{R}$ y $x_0 \in U$, entonces ella puede extenderse a una única solución maximal de \eqref{eq:517}.
	
	\noindent \textbf{Proposición 5.4.2} Sea $U \subseteq \mathbb{R}^n$ abierto y $X \in \mathcal{X}^0(U)$. Suponga que para todo $(t_0, x_0) \in \mathbb{R} \times U$, el P.V.I.
	\begin{equation}\label{eq:520}
		\left| 
		\begin{array}{lcl}
			x' &=& X(x) \\
			x(t_0) &=& x_0
		\end{array}
		\right.
	\end{equation}
	admite una única solución local definida en un intervalo abierto. Entonces se cumple:
	
	\begin{enumerate}
		\item Si $\psi_1 : J_1 \to U$ y $\psi_2 : J_2 \to U$ son soluciones de \eqref{eq:520} entonces $\left.\psi_1\right|_{J_1 \cap J_2} = \left.\psi_2\right|_{J_1 \cap J_2}$.
		\item El P.V.I. \eqref{eq:520} admite una única solución maximal $\varphi_M = \varphi_M(t_0, x_0)$ definida en el intervalo abierto $I_M = I_M(t_0, x_0)$.
	\end{enumerate}
	
	\noindent \textbf{Demostración.} 1.) Dados $t_0 \in \mathbb{R}$ y $x_0 \in U$, consideremos el conjunto
	\[
	C = \{t \in J_1 \cap J_2 \,;\, \psi_1(t) = \psi_2(t)\}.
	\]
	Es claro que $t_0 \in C$, luego $C \neq \emptyset$. También es claro que $C = (\psi_1 - \psi_2)^{-1}(\{0\}) \cap (J_1 \cap J_2)$, luego $C$ es cerrado en $J_1 \cap J_2$.
	
	Finalmente, sea $t' \in C$, entonces $\psi_1(t') = \psi_2(t') = x' \in U$. Consideremos el P.V.I.
	\begin{equation*}
		\left| 
		\begin{array}{lcl}
			x' &=& X(x) \\
			x(t') &=& x'
		\end{array}
		\right.
	\end{equation*}
	
	Por hipótesis este P.V.I. admite una única solución $\phi$ definida en un intervalo abierto $I' = I'(t', x')$, luego
	\[
	\left.\psi_1\right|_{I' \cap (J_1 \cap J_2)} = \left.\phi\right|_{I' \cap (J_1 \cap J_2)} = \left.\psi_2\right|_{I' \cap (J_1 \cap J_2)}.
	\]
	
	Se sigue que $t' \in I' \cap (J_1 \cap J_2) \subseteq C$, es decir $C$ es abierto.
	
	Desde que $J_1 \cap J_2$ es conexo, se sigue que $C = J_1 \cap J_2$.
	
	\noindent 2.) Dados $t_0 \in \mathbb{R}$ y $x_0 \in U$, consideremos el conjunto
	\[
	S = \{\psi : I_\psi \to \mathbb{R}^n \,;\, I_\psi \text{ es abierto y } \psi \text{ es solución del P.V.I. } \eqref{eq:520}\}.
	\]
	
	Por hipótesis $S \neq \emptyset$. Definamos $I_M = I_M(t_0, x_0) = \bigcup_{\psi \in S} I_\psi$, claramente $I_M$ es un intervalo abierto y $t_0 \in I_M$. Para definir $\varphi_M : I_M \to \mathbb{R}^n$ tomemos un $t \in I_M$, entonces existe $\psi \in S$ tal que $t \in I_\psi$. Definimos $\varphi_M(t) = \psi(t)$. Por la parte 1.), la función $\varphi_M$ está bien definida, y es claro que ella es la única solución maximal del P.V.I. \eqref{eq:520}. \hfill $\square$
	
	\noindent \textbf{Corolario.} Sea $U \subseteq \mathbb{R}^n$ un abierto y $X \in \mathcal{X}^1(U)$. Entonces para todo $(t_0, x_0) \in \mathbb{R} \times U$, el P.V.I.
	\begin{equation*}
		\left| 
		\begin{array}{lcl}
			x' &=& X(x) \\
			x(t_0) &=& x_0
		\end{array}
		\right.
	\end{equation*}
	admite una única solución maximal $\varphi_M = \varphi_M(t_0, x_0)$ definida en el intervalo abierto $I_M = I_M(t_0, x_0)$.
	
	\noindent \textbf{Notación:} Denotaremos respectivamente por $\omega_-(t_0, x_0)$ y $\omega_+(t_0, x_0)$ al extremo inferior y al superior del intervalo maximal $I_M(t_0, x_0)$, es decir
	\[
	I_M(t_0, x_0) = {]\omega_-(t_0, x_0), \omega_+(t_0, x_0)[}
	\]
	cuando no haya lugar a confusión con respecto a las condiciones iniciales, escribiremos simplemente $I_M = {]\omega_-, \omega_+[}$.
	
	\noindent \textbf{Observación:} No se puede garantizar que una solución $\varphi$ de un P.V.I. pueda extenderse a todo $\mathbb{R}$. En efecto, considere el P.V.I. escalar:
	\begin{equation*}
		\left| 
		\begin{array}{lcl}
			x' &=& 1 + x^2 \\
			x(0) &=& 0
		\end{array}
		\right.
	\end{equation*}
	cuya solución es dada por
	\[
	\begin{array}{rcccl}
		\varphi : & {]-\pi/2, \pi/2[} &\to& \mathbb{R} \\
		& t &\mapsto& \varphi(t) = \tan t
	\end{array}
	\]
	es claro que ella es una solución maximal del P.V.I. considerado.
	
	A continuación vamos a estudiar qué sucede con las soluciones maximales $\varphi_M : {]\omega_-, \omega_+[} \to U$, cuando $\omega_-$ u $\omega_+$ es un número real.
	
	\noindent \textbf{Lema 5.4.1} Sean $X \in \operatorname{Lip}(B_r[x_0], \mathbb{R}^n)$, $t_0 \in \mathbb{R}$, $M = \|X\|_{C(B_r[x_0])}$ y $\alpha = \frac{r}{M}$. Entonces para todo $(t_1, x_1) \in I_{\frac{\alpha}{2}}[t_0] \times B_{\frac{r}{2}}[x_0]$, el P.V.I.:
	\begin{equation}\label{eq:521}
		\left| 
		\begin{array}{lcl}
			x' &=& X(x) \\
			x(t_1) &=& x_1
		\end{array}
		\right.
	\end{equation}
	admite una única solución $\varphi_1$ (que depende de $(t_1, x_1)$) definida en el intervalo $I_{\frac{\alpha}{2}}[t_1] \subseteq I_\alpha[t_0]$, con valores en $B_{\frac{r}{2}}[x_1] \subseteq B_r[x_0]$.
	
	\noindent \textbf{Demostración.} Dados $t_1 \in I_{\frac{\alpha}{2}}[t_0]$ y $x_1 \in B_{\frac{r}{2}}[x_0]$, un fácil cálculo prueba que $I_{\frac{\alpha}{2}}[t_1] \subseteq I_\alpha[t_0]$ y $B_{\frac{r}{2}}[x_1] \subseteq B_r[x_0]$. Luego $M \ge \|X\|_{C(B_{\frac{r}{2}}[x_0])}$ y por el Teorema de Picard, existe una única solución $\varphi_1 = \varphi_1(t_1, x_1) : I_{\alpha_1}[t_1] \to B_{\frac{r}{2}}[x_0]$ del P.V.I. \eqref{eq:521}, en donde $\alpha_1 = \frac{r}{2M} = \frac{\alpha}{2}$. \hfill $\square$
	
	\noindent \textbf{Teorema 5.4.3} Sea $U \subseteq \mathbb{R}^n$ abierto, $(t_0, x_0) \in \mathbb{R} \times U$ y $X \in \mathcal{X}^1(U)$. Sea $\varphi$ solución maximal del P.V.I.:
	\begin{equation}\label{eq:522}
		\left| 
		\begin{array}{lcl}
			x' &=& X(x) \\
			x(t_0) &=& x_0
		\end{array}
		\right.
	\end{equation}
	definida en el intervalo maximal $I_M = {]\omega_-, \omega_+[}$. Si $\omega_+ \in \mathbb{R}$, (resp. $\omega_- \in \mathbb{R}$), entonces para todo $K \subseteq U$ compacto, existe $W_+ \subseteq \mathbb{R}$ vecindad abierta de $\omega_+$ (resp. $W_- \subseteq \mathbb{R}$ vecindad abierta de $\omega_-$), tal que $t \in W_+ \cap [t_0, \omega_+[$ (resp. $t \in W_- \cap {]\omega_-, t_0]}$) implica $\varphi(t) \notin K$.
	
	\noindent \textbf{Demostración.} Vamos a trabajar con $\omega_+$, en el otro caso la demostración es análoga. Procediendo por contradicción, supongamos que existe un $K \subseteq U$ subconjunto compacto tal que para toda vecindad abierta $W$ de $\omega_+$, existe $t = t_W \in W \cap [t_0, \omega_+[$ tal que $\varphi(t_W) \in K$.
	
	Dado $k \in \mathbb{Z}^^+$, consideremos el intervalo $W_k = I_{\frac{1}{k}}(\omega_+)$, el cual es una vecindad abierta de $\omega_+$. Por la hipótesis auxiliar, existe un $t_k \in W_k \cap [t_0, \omega_+[$ tal que $\varphi(t_k) \in K$. De esta manera, hemos construido una sucesión $(t_k) \subseteq [t_0, \omega_+[$, tal que $\lim_{k \to \infty} t_k = \omega_+$ y $(\varphi(t_k)) \subseteq K$.
	
	Como $K$ es compacto, existe una subsucesión $(\varphi(t_{j_k}))$ convergente de $(\varphi(t_k))$, es decir $\lim_{k \to \infty} \varphi(t_{j_k}) = x_0$ y $x_0 \in K$.
	
	Como $x_0 \in U$, existe un $r = r(x_0) > 0$ tal que $B_r[x_0] \subseteq U$ y desde que $X \in \mathcal{X}^1(U)$ entonces $X$ es Lipschitz en $B_r[x_0]$. Como por hipótesis $\omega_+ \in \mathbb{R}$, por el Teorema de Picard, existe una única solución $\psi$ del P.V.I.
	\begin{equation*}
		\left| 
		\begin{array}{lcl}
			x' &=& X(x) \\
			x(\omega_+) &=& x_0
		\end{array}
		\right.
	\end{equation*}
	definida en el intervalo $I_\alpha[\omega_+]$ donde $\alpha = \frac{r}{M}$ y $M = \|X\|_{C(B_r[x_0])}$.
	
	Del Lema 5.4.1 se tiene que para todo $(t_1, x_1) \in I_{\frac{\alpha}{2}}[\omega_+] \times B_{\frac{r}{2}}[x_0]$, el P.V.I.
	\begin{equation*}
		\left| 
		\begin{array}{lcl}
			x' &=& X(x) \\
			x(t_1) &=& x_1
		\end{array}
		\right.
	\end{equation*}
	admite una única solución definida en $I_{\frac{\alpha}{2}}[t_1]$ con valores en $B_{\frac{r}{2}}[x_1]$.
	
	Ahora bien, como $\lim_{k \to \infty} t_{j_k} = \omega_+$ y $\lim_{k \to \infty} \varphi(t_{j_k}) = x_0$, se sigue que existe un $k_0 \in \mathbb{Z}^+$ tal que $k \ge k_0$ implica
	\[
	\omega_+ - \frac{\alpha}{2} < t_{j_k} < \omega_+ \quad \text{y} \quad |\varphi(t_{j_k}) - x_0| < \frac{r}{2},
	\]
	luego, si fijamos $k \ge k_0$ tenemos que $(t_{j_k}, \varphi(t_{j_k})) \in I_{\frac{\alpha}{2}}[\omega_+] \times B_{\frac{r}{2}}[x_0]$, luego el P.V.I.:
	\begin{equation}\label{eq:523}
		\left| 
		\begin{array}{lcl}
			x' &=& X(x) \\
			x(t_{j_k}) &=& \varphi(t_{j_k})
		\end{array}
		\right.
	\end{equation}
	tiene una única solución $\psi_k : I_{\frac{\alpha}{2}}[t_{j_k}] \to B_{\frac{r}{2}}[\varphi(t_{j_k})] \subseteq B_r[x_0]$.
	
	Observe que $t_{j_k} < \omega_+ < t_{j_k} + \frac{\alpha}{2}$, de esta manera $\varphi : {]\omega_-, \omega_+[} \to U$
	también es solución de \eqref{eq:523}. Se sigue que $\varphi(t) = \psi_k(t)$, para todo $t \in {]\omega_-, \omega_+[} \cap I_{\frac{\alpha}{2}}[t_{j_k}]$.
	
	Definimos $\phi : {]\omega_-, t_{j_k} + \frac{\alpha}{2}[} \to \mathbb{R}^n$ como
	\[
	\left.\phi\right|_{]\omega_-, \omega_+[} = \varphi \quad \text{y} \quad \left.\phi\right|_{I_{\frac{\alpha}{2}}(t_{j_k})} = \psi_k
	\]
	
	Se sigue que $\phi$ es solución del P.V.I. \eqref{eq:522} cuyo dominio ${]\omega_-, t_{j_k} + \frac{\alpha}{2}[}$ contiene estrictamente al intervalo maximal ${]\omega_-, \omega_+[}$ lo cual constituye una contradicción. \hfill $\square$
	
	\section*{6.1 El Flujo Asociado a un Campo Vectorial}
	
	De ahora en adelante, sólo consideraremos P.V.I.'s con instante inicial $t_0 = 0$. Sea $U \subseteq \mathbb{R}^n$ abierto y $X \in \mathcal{X}^1(U)$. Dado $x_0 \in U$, sabemos que el P.V.I.
	\begin{equation}\label{eq:61}
		\left| 
		\begin{array}{lcl}
			x' &=& X(x) \\
			x(0) &=& x_0
		\end{array}
		\right.
	\end{equation}
	admite una única solución $\varphi : I_\alpha[0] \to B_r[x_0] \subseteq U$ (donde $\alpha = \frac{r}{M}$, siendo $M = \|X(x)\|_{C(B_r[x_0])}$), la cual puede ser extendida a su solución maximal $\varphi_{x_0} = \varphi_{(0,x_0)} : I(x_0) \to U$ (en donde $I(x_0) = {]\omega_-(x_0), \omega_+(x_0)[}$ es su intervalo maximal). Definamos el conjunto:
	\[
	D_X = \{(t,x) \,;\, x \in U \text{ y } t \in I(x)\} \subseteq \mathbb{R} \times \mathbb{R}^n
	\]
	y la función
	\[
	\begin{array}{rccc}
		\varphi_X : & D_X & \to & U \\
		& (t, x) & \mapsto & \varphi_X(t, x) = \varphi_x(t).
	\end{array}
	\]
	La función $\varphi_X$ es llamada flujo asociado al campo vectorial $X$.
	
	\noindent \textbf{Lema 6.1.1} Sean $U \subseteq \mathbb{R}^n$ abierto y $X \in \mathcal{X}^1(U)$. Si $y_0 = \varphi_{x_0}(t_0)$ con $t_0 \in I(x_0)$ entonces:
	\begin{enumerate}
		\item[\textnormal{i)}] $\varphi_{x_0}(0) = x_0$.
		\item[\textnormal{ii)}] $I(y_0) = I(x_0) - t_0 = \{t - t_0 \,;\, t \in I(x_0)\} = {]\omega_-(x_0) - t_0, \omega_+(x_0) - t_0[}$.
		\item[\textnormal{iii)}] $\varphi_{y_0}(t) = \varphi_{x_0}(t + t_0), \quad \forall \; t \in I(y_0)$.
	\end{enumerate}
	
	\noindent \textbf{Teorema 6.1.1} Sea $U \subseteq \mathbb{R}^n$ abierto y $X \in \mathcal{X}^1(U)$. Su flujo asociado $\varphi_X : D_X \to U$ satisface las siguientes propiedades:
	
	\begin{enumerate}
		\item $\varphi_X(0, x) = x, \quad \forall \; x \in U$.
		\item Si $(t, x) \in D_X$ y $(s, \varphi_x(t)) \in D_X$ entonces $(s + t, x) \in D_X$ y
		\[
		\varphi_X(s + t, x) = \varphi_X(s, \varphi_X(t, x)) \quad (\textit{Propiedad Fundamental del Flujo}).
		\]
	\end{enumerate}
	
	\noindent \textbf{Demostración.} 2.) Sea $y = \varphi_x(t)$. Por hipótesis y el lema anterior se tiene que $s \in I(y) = I(x) - t$ es decir $s + t \in I(x)$, por tanto $(s + t, x) \in D_X$. Además:
	\[
	\begin{aligned}
		\varphi_X(t + s, x) &= \varphi_x(t + s) = \varphi_y(s) = \varphi_X(s, y) = \varphi_X(s, \varphi_x(t)) \\
		&= \varphi_X(s, \varphi_X(t, x))
	\end{aligned}
	\]
	\hfill $\square$
	
	\textbf{Observación:} La aplicación $\varphi_t:U\to U$ definida por $\varphi_t(x) = \varphi_X(t,x)$ es un difeomorfismo de clase $C^k$ (si es que $X$ es de clase $C^k$) y su inversa es $\varphi_t^{-1} = \varphi_{-t}$.
	
	\section*{Definición del Corchete de Lie mediante Pull-back}
	
	Existe una caracterización geométrica fundamental que define el corchete de Lie midiendo cómo varía el campo $Y$ al ser transportado de regreso ("pull-back") a lo largo de las curvas integrales del campo $X$.
	
	Sean $U \subseteq \mathbb{R}^n$ un abierto, $X, Y \in \mathcal{X}^1(U)$ dos campos vectoriales y sea $\varphi_t$ el flujo local asociado al campo $X$.
	
	\begin{definicion}[Pull-back de un Campo Vectorial por el Flujo]
		Dado que para cada $t$ suficientemente pequeño $\varphi_t : U \to U$ es un difeomorfismo local, el \textbf{pull-back} de $Y$ a lo largo de $\varphi_t$, denotado por $\varphi_t^* Y$, es el campo vectorial definido en $p \in U$ mediante:
		\[
		(\varphi_t^* Y)(p) = \left( d\varphi_t(p) \right)^{-1} Y(\varphi_t(p)) = d\varphi_{-t}(\varphi_t(p)) \cdot Y(\varphi_t(p))
		\]
		donde $d\varphi_t(p)$ denota la diferencial (matriz jacobiana) de la aplicación $\varphi_t$ evaluada en $p$.
	\end{definicion}
	
	\begin{definicion}[Corchete de Lie como Derivada de Lie]
		El \textbf{corchete de Lie} $[X, Y]$ se define como la \textbf{derivada de Lie} del campo $Y$ a lo largo de $X$ (denotada $\mathcal{L}_X Y$), la cual se obtiene como la derivada en $t=0$ del pull-back de $Y$:
		\[
		[X, Y](p) \equiv (\mathcal{L}_X Y)(p) = \lim_{t \to 0} \frac{(\varphi_t^* Y)(p) - Y(p)}{t} = \left. \frac{d}{dt} \right|_{t=0} (\varphi_t^* Y)(p)
		\]
	\end{definicion}
	
	\subsection*{Interpretación Geométrica y Conmutatividad}
	
	\begin{itemize}
		\item \textbf{Significado físico/geométrico:} La cantidad $(\varphi_t^* Y)(p)$ toma el vector $Y$ en el punto trasladado $\varphi_t(p)$ y lo retrocede al punto original $p$ usando la diferencial inversa $d\varphi_{-t}$. La derivada en $t=0$ mide qué tanto falla el campo $Y$ en ser invariante bajo la acción del flujo de $X$.
		\item \textbf{Conmutatividad de flujos:} Si $\psi_s$ denota el flujo del campo $Y$, se verifica el siguiente resultado fundamental:
		\[
		[X, Y] = 0 \iff \varphi_t \circ \psi_s = \psi_s \circ \varphi_t \quad \text{(los flujos conmutan)}.
		\]
	\end{itemize}
	
	\begin{teorema}
		Sean $X, Y \in \mathcal{X}^1(U)$ dos campos vectoriales de clase $C^1$, y sean $\varphi_t$ y $\psi_s$ sus respectivos flujos locales. Entonces:
		\[
		[X, Y] = 0 \iff \varphi_t \circ \psi_s = \psi_s \circ \varphi_t
		\]
		(para todo $t, s$ suficientemente pequeños).
	\end{teorema}
	
	\begin{proof}
		\textbf{$(\Leftarrow)$:} Supongamos que $\varphi_t(\psi_s(p)) = \psi_s(\varphi_t(p))$. Derivando con respecto a $s$ en $s = 0$:
		\[
		d\varphi_t(p) \cdot Y(p) = Y(\varphi_t(p)) \implies (\varphi_t^* Y)(p) = Y(p)
		\]
		Derivando esta última expresión con respecto a $t$ en $t = 0$ se obtiene $[X, Y](p) = 0$.
		
		\textbf{$(\Rightarrow)$:} Supongamos que $[X, Y] = 0$, lo que implica que el \textit{pull-back} es constante ($\varphi_t^* Y = Y$). Para un punto $p$ fijo, consideremos las curvas $f(t) = \varphi_t(\psi_s(p))$ y $g(t) = \psi_s(\varphi_t(p))$. Ambas parten de $\psi_s(p)$ en $t=0$ y satisfacen la misma EDO $\frac{dx}{dt} = X(x)$ (la segunda utilizando la invarianza del campo $X$ bajo el flujo de $Y$). Por la unicidad de las soluciones del P.V.I., se concluye que $f(t) = g(t)$, es decir, $\varphi_t \circ \psi_s = \psi_s \circ \varphi_t$.
	\end{proof}
	
	\section*{Corchete de Lie de Campos Vectoriales}
	
	Sean $U \subseteq \mathbb{R}^n$ un conjunto abierto y $X, Y \in \mathcal{X}^1(U)$ dos campos vectoriales de clase $C^1$ definidos sobre $U$.
	
	\begin{definicion}[Corchete de Lie]
		El \textbf{corchete de Lie} (o paréntesis de Lie) de los campos vectoriales $X$ e $Y$, denotado por $[X, Y]$, es el campo vectorial definido sobre $U$ dado por:
		\[
		[X, Y](p) = D Y(p) \cdot X(p) - D X(p) \cdot Y(p), \quad \forall \; p \in U
		\]
		donde $D X(p)$ y $D Y(p)$ denotan la matriz jacobiana (o diferencial) de los campos $X$ e $Y$ evaluadas en el punto $p$, respectivamente.
		
		En componentes, si $X = (X_1, \dots, X_n)$ e $Y = (Y_1, \dots, Y_n)$, la $i$-ésima componente de $[X, Y]$ está dada por:
		\[
		[X, Y]_i = \sum_{j=1}^n \left( X_j \frac{\partial Y_i}{\partial x_j} - Y_j \frac{\partial X_i}{\partial x_j} \right), \quad i = 1, \dots, n
		\]
	\end{definicion}
	
	\subsection*{Caracterización como Operador Diferencial}
	
	Si se consideran los campos vectoriales $X, Y$ como derivaciones sobre el álgebra de funciones $C^\infty(U)$, el corchete de Lie es el único campo vectorial que satisface:
	\[
	[X, Y](f) = X(Y(f)) - Y(X(f)), \quad \forall \; f \in C^\infty(U)
	\]
	donde $X(f) = \nabla f \cdot X = \sum_{i=1}^n X_i \frac{\partial f}{\partial x_i}$ denota la derivada direccional de $f$ a lo largo de $X$.
	
	\subsection*{Propiedades Fundamentales}
	
	Para cualesquiera $X, Y, Z \in \mathcal{X}^1(U)$, $a, b \in \mathbb{R}$ y $f, g \in C^\infty(U)$:
	\begin{enumerate}
		\item \textbf{Bilinealidad:} $[a X + b Y, Z] = a [X, Z] + b [Y, Z]$.
		\item \textbf{Antisimetría (o conmutatividad):} $[X, Y] = -[Y, X]$.
		\item \textbf{Identidad de Jacobi:} 
		\[
		[X, [Y, Z]] + [Y, [Z, X]] + [Z, [X, Y]] = 0
		\]
		\item \textbf{Regla de Leibniz:} 
		\[
		[X, f Y] = f [X, Y] + X(f) Y
		\]
	\end{enumerate}
	
	\subsection*{Ejercicio 1: Demostración de que $X$ es una derivación}
	
	Sean $f, g \in C^\infty(U)$, $c \in \mathbb{R}$ y $X \in \mathcal{X}^\infty(U)$ un campo vectorial. La aplicación $X : C^\infty(U) \to C^\infty(U)$ definida por $X(f)(p) = \langle \nabla f(p), X(p) \rangle$ cumple:
	
	\subsubsection*{1. Linealidad}
	Para cualquier $p \in U$:
	\[
	\begin{aligned}
		X(f + c g)(p) &= \langle \nabla(f + cg)(p), X(p) \rangle \\
		&= \langle \nabla f(p) + c \nabla g(p), X(p) \rangle \\
		&= \langle \nabla f(p), X(p) \rangle + c \langle \nabla g(p), X(p) \rangle \\
		&= X(f)(p) + c X(g)(p)
	\end{aligned}
	\]
	Por lo tanto, $X(f + cg) = X(f) + c X(g)$.
	
	\subsubsection*{2. Regla de Leibniz (Regla del Producto)}
	Por la regla del producto para el gradiente en cálculo multivariable, sabemos que:
	\[
	\nabla(fg)(p) = f(p) \nabla g(p) + g(p) \nabla f(p)
	\]
	Evaluando $X(fg)$ en el punto $p$:
	\[
	\begin{aligned}
		X(fg)(p) &= \langle \nabla(fg)(p), X(p) \rangle \\
		&= \langle f(p) \nabla g(p) + g(p) \nabla f(p), X(p) \rangle \\
		&= f(p) \langle \nabla g(p), X(p) \rangle + g(p) \langle \nabla f(p), X(p) \rangle \\
		&= f(p) X(g)(p) + g(p) X(f)(p)
	\end{aligned}
	\]
	Por lo tanto, en todo $U$ se satisface la identidad:
	\[
	X(fg) = f X(g) + g X(f)
	\]
	lo que prueba que el campo vectorial $X$ actúa como una \textbf{derivación} sobre el álgebra $C^\infty(U)$.
	
	\section*{Demostración de la Equivalencia de las Definiciones del Corchete de Lie}
	
	\begin{teorema}
		Sean $U \subseteq \mathbb{R}^n$ un abierto, $X, Y \in \mathcal{X}^1(U)$ campos vectoriales, y sea $\varphi_t$ el flujo local generado por el campo $X$. La derivada de Lie del campo $Y$ a lo largo del flujo de $X$ coincide con la expresión analítica del corchete de Lie:
		\[
		\left. \frac{d}{dt} \right|_{t=0} (\varphi_t^* Y)(p) = DY(p) X(p) - DX(p) Y(p) \equiv [X, Y](p).
		\]
	\end{teorema}
	
	\begin{proof}
		Sea $p \in U$. Recordemos que el \textit{pull-back} de $Y$ por la aplicación $\varphi_t$ en el punto $p$ está dado por:
		\[
		(\varphi_t^* Y)(p) = \left( d\varphi_t(p) \right)^{-1} Y(\varphi_t(p))
		\]
		donde $d\varphi_t(p) = D_x \varphi(t, p)$ representa la matriz jacobiana de la aplicación espacial $x \mapsto \varphi(t, x)$ evaluada en $p$.
		
		Para calcular la derivada con respecto a $t$ en $t=0$, desarrollaremos en serie de Taylor alrededor de $t=0$ los dos factores que componen el \textit{pull-back}:
		
		\subsubsection*{Paso 1: Expansión de $Y(\varphi_t(p))$}
		Por la definición de flujo, $\varphi(0, p) = p$ y $\left.\frac{\partial \varphi}{\partial t}\right|_{t=0} = X(p)$. Aplicando la regla de la cadena para la expansión de Taylor de $Y$ en $t=0$:
		\[
		Y(\varphi_t(p)) = Y(p) + t \, DY(p) \cdot X(p) + O(t^2)
		\]
		
		\subsubsection*{Paso 2: Expansión del operador diferencial inverso $(d\varphi_t(p))^{-1}$}
		Por el Teorema de Schwarz sobre la conmutación de derivadas parciales, la derivada temporal de la diferencial espacial satisface:
		\[
		\left. \frac{d}{dt} \right|_{t=0} d\varphi_t(p) = d \left( \left. \frac{d\varphi_t}{dt} \right|_{t=0} \right)(p) = dX(p) = DX(p)
		\]
		Como $d\varphi_0(p) = I$ (la matriz identidad), la expansión de Taylor de la matriz $d\varphi_t(p)$ es:
		\[
		d\varphi_t(p) = I + t \, DX(p) + O(t^2)
		\]
		Utilizando la serie de Neumann para la inversa de matrices $(I + t A)^{-1} = I - t A + O(t^2)$, obtenemos:
		\[
		\left( d\varphi_t(p) \right)^{-1} = I - t \, DX(p) + O(t^2)
		\]
		
		\subsubsection*{Paso 3: Producto de las expansiones}
		Sustituyendo ambas expansiones en la expresión del \textit{pull-back}:
		\[
		\begin{aligned}
			(\varphi_t^* Y)(p) &= \left( I - t \, DX(p) + O(t^2) \right) \left( Y(p) + t \, DY(p) X(p) + O(t^2) \right) \\
			&= Y(p) + t \, DY(p) X(p) - t \, DX(p) Y(p) + O(t^2) \\
			&= Y(p) + t \Big( DY(p) X(p) - DX(p) Y(p) \Big) + O(t^2)
		\end{aligned}
		\]
		
		\subsubsection*{Paso 4: Cálculo del límite (derivada en $t=0$)}
		Finalmente, calculando la derivada por cociente incremental en $t=0$:
		\[
		\begin{aligned}
			\left. \frac{d}{dt} \right|_{t=0} (\varphi_t^* Y)(p) &= \lim_{t \to 0} \frac{(\varphi_t^* Y)(p) - Y(p)}{t} \\
			&= \lim_{t \to 0} \frac{t \Big( DY(p) X(p) - DX(p) Y(p) \Big) + O(t^2)}{t} \\
			&= DY(p) X(p) - DX(p) Y(p)
		\end{aligned}
		\]
		Esto demuestra que ambas definiciones representan exactamente el mismo campo vectorial.
	\end{proof}
	
	%================================================
	% CAPÍTULO 7
	%================================================
	\section*{Capítulo 7: Introducción a la Teoría Geométrica Local}
	
	En el presente capítulo iniciamos el estudio del comportamiento geométrico de la foliación generada por un campo vectorial en la vecindad de un punto $p$ dado. Veremos que si $p$ es un punto regular, el problema queda resuelto satisfactoriamente, pero si $p$ es singular surge una teoría interesante de la cual sólo haremos una introducción.
	
	\subsection*{7.1 Estructura local de los puntos regulares}
	
	En esta sección, probaremos que si $p$ es un punto regular de un campo $X$, entonces en una vecindad suficientemente pequeña de $p$ podemos ``enderezar'' las órbitas de $X$. Para su demostración usaremos algunos resultados previos.
	
	En lo que sigue, el símbolo $\langle v_1, v_2, \dots, v_m \rangle$ denota al subespacio vectorial de $\mathbb{R}^n$ generado por los vectores $v_1, v_2, \dots, v_m$.
	
	\noindent \textbf{Definición 7.1.1} Sea $U \subseteq \mathbb{R}^n$ un abierto, $X \in \mathcal{X}^k(U)$ y $p \in U$. Una \textit{sección transversal local de clase $C^r$} al campo $X$ en el punto $p$, es una función $\sigma : V^{n-1} \to U$ de clase $C^r$ (en donde $V^{n-1}$ es un abierto conexo de $\mathbb{R}^{n-1}$) que satisface las siguientes propiedades:
	\begin{enumerate}[label=\roman*)]
		\item $\sigma$ es una inmersión de $V^{n-1}$ en $U$ (es decir, $\sigma'(x) \in \mathcal{L}(\mathbb{R}^{n-1}, \mathbb{R}^n)$ es inyectiva, $\forall \; x \in V^{n-1}$).
		\item $\sigma : V^{n-1} \to \sigma(V^{n-1}) = \Sigma$ es un homeomorfismo.
		\item $p \in \Sigma$.
		\item $\sigma'(x)[\mathbb{R}^{n-1}_x] \oplus \langle X(\sigma(x)) \rangle = \mathbb{R}^n_{\sigma(x)}, \quad \forall \; x \in V^{n-1}$.
	\end{enumerate}
	
	\noindent \textbf{Ejemplo 7.1.1} Sea $X \in \mathcal{X}^\infty(\mathbb{R}^2)$ definido por
	\[
	\begin{array}{rccc}
		X : & \mathbb{R}^2 & \to & \mathbb{R}^2 \\
		& (x, y) & \mapsto & X(x, y) = (x, -y).
	\end{array}
	\]
	Sea $p = (1, 0) \in \mathbb{R}^2$ y consideremos la función
	\[
	\begin{array}{rccc}
		\sigma : & {]-1, 1[} & \to & \mathbb{R}^2 \\
		& t & \mapsto & \sigma(t) = (1, -t).
	\end{array}
	\]
	Afirmo que $\sigma$ es una sección transversal local de clase $C^\infty$ al campo $X$ en el punto $p$. En efecto:
	\begin{enumerate}[label=\roman*)]
		\item $\sigma'(t) = (0, -1) \neq (0, 0), \quad \forall \; t \in {]-1, 1[}$. Así, $\sigma$ es una inmersión.
		\item $\Sigma = \sigma({]-1, 1[}) = \{(1, t) \,;\, t \in {]-1, 1[}\}$, luego $\sigma : {]-1, 1[} \to \Sigma$ es un homeomorfismo.
		\item $p = \sigma(0) \in \Sigma$.
		\item $\sigma'(t)[\mathbb{R}] = \langle (0, 1) \rangle$ y $X(\sigma(t)) = X(1, -t) = (1, t), \quad \forall \; t \in {]-1, 1[}$. Luego:
		\[
		\sigma'(t)[\mathbb{R}] \oplus \langle X(\sigma(t)) \rangle = \mathbb{R}^2, \quad \forall \; t \in {]-1, 1[}.
		\]
	\end{enumerate}
	
	El siguiente resultado prueba la existencia de secciones transversales en vecindades de puntos regulares de campos de clase $C^k$.
	
	\noindent \textbf{Proposición 7.1.1} \textit{Si $U \subseteq \mathbb{R}^n$ es un abierto, $X \in \mathcal{X}^k(U)$ y $p \in U$ es un punto regular de $X$, entonces existe una sección transversal local $\sigma$ de clase $C^\infty$ al campo $X$ en el punto $p$ tal que $\sigma(0) = p$.}
	
	\begin{proof}
		Como $p$ es un punto regular de $X$, entonces $X(p) \neq 0$, luego existen vectores $v_1, \dots, v_{n-1} \in \mathbb{R}^n$ tales que $\{v_1, \dots, v_{n-1}, X(p)\}$ es una base de $\mathbb{R}^n$.
		
		Afirmo que existe un $\delta > 0$ tal que $\{v_1, \dots, v_{n-1}, X(q)\}$ es base de $\mathbb{R}^n$, $\forall \; q \in B_\delta(p) \subseteq U$. En efecto, puesto que en caso contrario se tendría que para cualquier $\delta > 0$ existiría un $q_\delta \in B_\delta(p)$ tal que $X(q_\delta) \in \langle v_1, v_2, \dots, v_{n-1} \rangle$.
		
		De esta manera, dado cualquier $m \in \mathbb{Z}^+$, existe un punto $q_m \in B_{1/m}(p)$ tal que
		\[
		X(q_m) \in \langle v_1, v_2, \dots, v_{n-1} \rangle.
		\]
		Hemos construido así una sucesión $(q_m)_{m \in \mathbb{Z}^+} \subseteq U$ tal que $\lim_{m \to \infty} q_m = p$ y $X(q_m) \in \langle v_1, v_2, \dots, v_{n-1} \rangle$. Luego:
		\[
		\lim_{m \to \infty} X(q_m) = X(p) \quad \text{y} \quad X(p) \in \overline{\langle v_1, v_2, \dots, v_{n-1} \rangle} = \langle v_1, v_2, \dots, v_{n-1} \rangle,
		\]
		lo cual es una contradicción.
		
		Sea $M = \max\{\|v_1\|, \|v_2\|, \dots, \|v_{n-1}\|\}$ y considero $\varepsilon < \frac{\delta}{M(n-1)}$. Defino
		\[
		V^{n-1} = \{x \in \mathbb{R}^{n-1} \,;\, \|x\| < \varepsilon\}
		\]
		y
		\[
		\begin{array}{rccc}
			\sigma : & V^{n-1} & \to & \mathbb{R}^n \\
			& x & \mapsto & \sigma(x) = p + \displaystyle\sum_{j=1}^{n-1} x_j v_j.
		\end{array}
		\]
		Observe que
		\[
		\|\sigma(x) - p\| = \left\| \sum_{j=1}^{n-1} x_j v_j \right\| \le \sum_{j=1}^{n-1} |x_j| \, \|v_j\| \le M \sum_{j=1}^{n-1} |x_j| \le M(n-1)\|x\| < M(n-1)\varepsilon < \delta,
		\]
		es decir, $\sigma(x) \in B_\delta(p) \subseteq U, \quad \forall \; x \in V^{n-1}$. Afirmo que $\sigma : V^{n-1} \to U$ es la sección transversal buscada. En efecto, claramente $\sigma$ es de clase $C^\infty$ y un fácil cálculo muestra que, dado $a = (a_1, \dots, a_{n-1}) \in \mathbb{R}^{n-1}$, se tiene
		\[
		\sigma'(x)(a) = \sum_{j=1}^{n-1} a_j v_j,
		\]
		luego se cumple:
		\begin{enumerate}[label=\roman*)]
			\item $\sigma'(x)(a) = 0 \iff \sum_{j=1}^{n-1} a_j v_j = 0 \iff a_j = 0, \; \forall \; 1 \le j \le n-1 \iff a = 0$. Por lo tanto, $\sigma'(x)$ es inyectiva, $\forall \; x \in V^{n-1}$.
			\item Si $\Sigma = \sigma(V^{n-1})$ entonces $\sigma : V^{n-1} \to \Sigma$ es un homeomorfismo.
			\item $p = \sigma(0) \in \Sigma$.
			\item Como $\sigma'(x)(e_i) = v_i$ para todo $i = 1, \dots, n-1$, se tiene que
			\[
			\sigma'(x)[\mathbb{R}^{n-1}_x] = \langle v_1, v_2, \dots, v_{n-1} \rangle, \quad \forall \; x \in V^{n-1}.
			\]
			Se sigue que
			\[
			\sigma'(x)[\mathbb{R}^{n-1}_x] \oplus \langle X(\sigma(x)) \rangle = \mathbb{R}^n_{\sigma(x)}, \quad \forall \; x \in V^{n-1}. \qedhere
			\]
		\end{enumerate}
	\end{proof}
	
	Otro resultado que necesitamos es el siguiente:
	
	\noindent \textbf{Teorema 7.1.2 (Teorema de la Función Inversa)} \textit{Sea $U \subseteq \mathbb{R}^n$ abierto, $f : U \to \mathbb{R}^n$ función de clase $C^k$ tal que $f'(a) \in \operatorname{GL}(\mathbb{R}^n)$ (donde $a \in U$), entonces existen $V_a \subseteq U$ y $W_a \subseteq \mathbb{R}^n$ abiertos con $a \in V_a$, $f(a) \in W_a$, tales que $\restr{f}{V_a} : V_a \to W_a$ es un difeomorfismo de clase $C^k$.}
	
	La prueba del Teorema de la Función Inversa el lector la puede encontrar en [16].
	
	\noindent \textbf{Teorema 7.1.3 (Teorema del Flujo Tubular)} \textit{Sea $U \subseteq \mathbb{R}^n$ abierto, $X \in \mathcal{X}^k(U)$ y $p \in U$ un punto regular. Dada una sección transversal local $\sigma : V \to \Sigma$ de clase $C^\infty$ a $X$ en $p$, con $0 \in V$ y $\sigma(0) = p$, entonces existe una vecindad $W_p$ de $p$ en $U$ y un difeomorfismo de clase $C^k$}
	\[
	h : W_p \to I_\varepsilon(0) \times B
	\]
	\textit{en donde $\varepsilon > 0$ y $B$ es una bola abierta en $\mathbb{R}^{n-1}$ centrada en el origen, tales que:}
	\begin{enumerate}[label=\roman*)]
		\item $h(\Sigma \cap W_p) = \{0\} \times B$.
		\item $h$ \textit{es una conjugación de clase $C^k$ entre $\restr{X}{W_p}$ y el campo constante}
		\[
		Y : I_\varepsilon(0) \times B \to \mathbb{R}^n, \quad Y = (1, 0, 0, \dots, 0).
		\]
	\end{enumerate}
	
	\begin{proof}
		Sea $\varphi : D \to U$ el flujo de $X$ y denotemos
		\[
		\widetilde{D} = \{(t, u) \in \mathbb{R} \times V \,;\, (t, \sigma(u)) \in D\} \subseteq \mathbb{R} \times \mathbb{R}^{n-1}.
		\]
		Definimos la función
		\[
		\begin{array}{rccc}
			F : & \widetilde{D} & \to & U \\
			& (t, u) & \mapsto & F(t, u) = \varphi(t, \sigma(u)).
		\end{array}
		\]
		Claramente $F$ es de clase $C^k$. Geométricamente la función $F$ transforma líneas horizontales (órbitas del campo $Y$) en órbitas de $X$ que en el instante $0$ pasan por la sección transversal.
		
		Observe que $F(0, \mathbf{0}) = p$, vamos a probar que $F$ es un difeomorfismo local en $(0, \mathbf{0}) \in \mathbb{R} \times \mathbb{R}^{n-1}$. Por el Teorema de la Función Inversa, es suficiente probar que $F'(0, \mathbf{0})$ es un isomorfismo. En primer lugar, usando las notaciones del Análisis en varias variables reales tenemos
		\begin{equation}\label{eq:71}
			F'(0, \mathbf{0}) = \left[ \frac{\partial F}{\partial t}(0, \mathbf{0}), \frac{\partial F}{\partial u_1}(0, \mathbf{0}), \dots, \frac{\partial F}{\partial u_{n-1}}(0, \mathbf{0}) \right] \in \mathbb{R}^{n \times n}.
		\end{equation}
		Ahora bien, por un lado tenemos
		\[
		\frac{\partial F}{\partial t}(0, \mathbf{0}) = \frac{\partial \varphi}{\partial t}(0, \sigma(\mathbf{0})) = \frac{\partial \varphi}{\partial t}(0, p) = X(\varphi(0, p)) = X(p)
		\]
		y por otro lado
		\[
		\frac{\partial F}{\partial u_j}(0, \mathbf{0}) = \frac{\partial \varphi}{\partial t}(0, p) \cdot 0 + \frac{\partial \varphi}{\partial x_1}(0, p) \frac{\partial \sigma_1}{\partial u_j}(0) + \frac{\partial \varphi}{\partial x_2}(0, p) \frac{\partial \sigma_2}{\partial u_j}(0) + \dots + \frac{\partial \varphi}{\partial x_n}(0, p) \frac{\partial \sigma_n}{\partial u_j}(0).
		\]
		Pero $\varphi(0, x) = x$ implica $\frac{\partial \varphi}{\partial x_i}(0, p) = e_i$ para todo $1 \le i \le n$, de donde:
		\[
		\begin{aligned}
			\frac{\partial F}{\partial u_j}(0, \mathbf{0}) &= \frac{\partial \sigma_1}{\partial u_j}(0) e_1 + \frac{\partial \sigma_2}{\partial u_j}(0) e_2 + \dots + \frac{\partial \sigma_n}{\partial u_j}(0) e_n \\
			&= \left( \frac{\partial \sigma_1}{\partial u_j}(0), \frac{\partial \sigma_2}{\partial u_j}(0), \dots, \frac{\partial \sigma_n}{\partial u_j}(0) \right) \\
			&= \frac{\partial \sigma}{\partial u_j}(0) = \sigma'(\mathbf{0}) e_j, \quad \forall \; 1 \le j \le n-1.
		\end{aligned}
		\]
		Luego, reemplazando en \eqref{eq:71}:
		\[
		F'(0, \mathbf{0}) = \left[ X(p), \sigma'(0)e_1, \dots, \sigma'(0)e_{n-1} \right].
		\]
		Como $\sigma = (\sigma_1, \sigma_2, \dots, \sigma_n) : V \to \Sigma$ es una sección transversal, se cumple:
		\[
		\sigma'(0)[\mathbb{R}^{n-1}] \oplus \langle X(p) \rangle = \mathbb{R}^n,
		\]
		luego $F'(0, \mathbf{0}) \in \operatorname{GL}(\mathbb{R}^n)$. Por el Teorema de la Función Inversa, existe $\varepsilon = \varepsilon(p) > 0$, existe $B \subseteq \mathbb{R}^{n-1}$ bola abierta centrada en $0 \in \mathbb{R}^{n-1}$ y existe $W_p$ vecindad abierta de $p$ en $U$ tal que
		\[
		\restr{F}{I_\varepsilon(0) \times B} : I_\varepsilon(0) \times B \to W_p
		\]
		es un difeomorfismo de clase $C^k$. Consideremos
		\[
		h = \left( \restr{F}{I_\varepsilon(0) \times B} \right)^{-1} : W_p \to I_\varepsilon(0) \times B.
		\]
		Como $(0, u) \in I_\varepsilon(0) \times B$ es tal que $F(0, u) = \varphi(0, \sigma(u)) = \sigma(u) \in \Sigma \cap W_p$, concluimos que $h(\Sigma \cap W_p) = \{0\} \times B$.
		
		Finalmente, para $(t, u) \in I_\varepsilon(0) \times B$, tenemos
		\[
		\begin{aligned}
			(h^{-1})'(t, u) \cdot Y(t, u) &= F'(t, u) \cdot (1, 0, \dots, 0) = \frac{\partial F}{\partial t}(t, u) \\
			&= \frac{\partial \varphi}{\partial t}(t, \sigma(u)) = X(\varphi(t, \sigma(u))) \\
			&= X(F(t, u)) = X(h^{-1}(t, u)),
		\end{aligned}
		\]
		luego $h^{-1}$ es una $C^k$-conjugación entre $Y$ y $X$, y por lo tanto $h$ es una $C^k$-conjugación entre $\restr{X}{W_p}$ e $Y$.
	\end{proof}
	
	En muchas ocasiones estamos interesados en conocer cómo son las órbitas de un campo en la vecindad de un punto dado. Surge de esta manera el concepto de \textit{conjugación local}.
	
	\noindent \textbf{Definición 7.1.2} Sean $U_1, U_2 \subseteq \mathbb{R}^n$ abiertos, $X_1 \in \mathcal{X}^k(U_1)$, $X_2 \in \mathcal{X}^k(U_2)$, $p_1 \in U_1$, $p_2 \in U_2$ y $\varphi_1 : D_1 \to U_1$, $\varphi_2 : D_2 \to U_2$ los flujos asociados a los campos $X_1$ y $X_2$ respectivamente. Decimos que $X_1$ y $X_2$ son \textbf{localmente topológicamente} (resp. $C^r$) \textbf{conjugados} en $p_1$ y $p_2$ si y sólo si existen vecindades abiertas $V_1 \subseteq U_1$ y $V_2 \subseteq U_2$ de $p_1$ y $p_2$ respectivamente y un homeomorfismo (resp. difeomorfismo de clase $C^r$) $h : V_1 \to V_2$ llamado \textbf{conjugación topológica} (resp. de clase $C^r$) \textbf{local} tal que
	\[
	h(\varphi_1(t, x)) = \varphi_2(t, h(x)), \quad \forall \; x \in V_1, \quad \forall \; t \in I(x).
	\]
\end{document}